% ---
% Reemplaza líneas 1347-1390 de Docs/Technical_Report.tex
% Decisiones: D-031, H-012, H-034
% ---

\subsubsection{\textit{PauliFeatureMap}: generalización por elección de los operadores de Pauli}

El circuito \textit{PauliFeatureMap} constituye la generalización del esquema de
Havl{\'{\i}}{\v{c}}ek et al. que permite seleccionar subconjuntos arbitrarios de
operadores de Pauli ($X$, $Y$, $Z$) en el bloque unitario de codificación. El operador
adopta la forma general:
\begin{equation}
  U_{\text{Pauli}}(\mathbf{x}) = \left[\mathcal{U}_\Phi^{(P)}(\mathbf{x})\,H^{\otimes n}\right]^r,
  \label{eq:paulifeaturemap}
\end{equation}
donde el generador diagonal o no diagonal se define a partir de una familia de cadenas
de Pauli $\{P_S\}$:
\begin{equation}
  \mathcal{U}_\Phi^{(P)}(\mathbf{x}) = \prod_{P \in \mathcal{P}}
  \exp\!\left(-i \sum_{S} \phi_S(\mathbf{x})\,\bigotimes_{j \in S} P_j\right),
  \label{eq:pauli_expansion}
\end{equation}
donde $\mathcal{P}$ es la lista de cadenas de Pauli elegida, en el orden en que se aplican. Si los
términos no conmutan, como $Y$ y $ZZ$, el bloque es un producto de exponenciales y no una
sola exponencial: en C5 se aplica primero $\exp\bigl(-i\sum_i x_i Y_i\bigr)$ y después
$\exp\bigl(-i\sum_{i<j}(\pi-x_i)(\pi-x_j)Z_iZ_j\bigr)$.
En el presente trabajo, la condición experimental C5 se define formalmente como
\texttt{pauli\_feature\_map(paulis=['Y', 'ZZ'], reps=1)}, manteniendo el mismo número
de qubits ($k = 12$) e idéntica conectividad de entrelazamiento de segundo orden
($ZZ$) que el mapa C4 (\texttt{zz\_feature\_map}). La
comparación entre C4 y C5 aísla con precisión el efecto de sustituir el eje de rotación
de la codificación de primer orden (eje $Z$ frente a eje $Y$), preservando
inalterada la estructura del entrelazamiento de dos cuerpos.

La elección de operadores de Pauli en la codificación de un solo qubit tiene consecuencias
geométricas sobre la esfera de Bloch:
\begin{enumerate}
  \item \textbf{El término en $X$ no codifica nada:} Tras la capa inicial de Hadamard,
        el estado del registro se encuentra en $|+\rangle^{\otimes n} = \left(\frac{|0\rangle + |1\rangle}{\sqrt{2}}\right)^{\otimes n}$.
        Dado que el estado $|+\rangle$ es un autoestado del operador de Pauli $X$ con
        autovalor $+1$ ($X|+\rangle = |+\rangle$), la aplicación del operador unitario de
        evolución bajo $X$ satisface:
        \begin{equation}
          e^{-i\,\phi(x_i)\,X}|+\rangle = e^{-i\,\phi(x_i)}\,|+\rangle.
          \label{eq:x_eigenstate_phase}
        \end{equation}
        La variable $x_i$ induce únicamente una fase global no física sobre el qubit, sin
        alterar las amplitudes relativas ni modificar la densidad de probabilidad en
        ninguna base. Un mapa basado en $\{X, ZZ\}$ anula la codificación
        de primer orden y degenera en un circuito de puro entrelazamiento $ZZ$.
  \item \textbf{Rotación a lo largo del ecuador (Pauli $Z$, C4):} El generador $Z$ desplaza
        el estado $|+\rangle$ a lo largo del plano ecuatorial de la esfera de Bloch, alterando
        la fase relativa $\varphi$ entre $|0\rangle$ y $|1\rangle$ sin cambiar las
        probabilidades marginales de ocupación de los polos.
  \item \textbf{Rotación fuera del ecuador hacia los polos (Pauli $Y$, C5):} El generador
        $Y$ actúa perpendicularmente al estado ecuatorial $|+\rangle$, rotándolo a lo largo
        de un meridiano hacia los polos $|0\rangle$ y $|1\rangle$. Esta transformación
        modifica de forma directa las probabilidades de medición en la base computacional
        antes de la aplicación del entrelazamiento.
\end{enumerate}

Qiskit construye el término $\exp(-i x_i Y)$ como $\sqrt{X}\,P(2x_i)\,\sqrt{X}^\dagger$, que equivale a $R_Y(2x_i)$ salvo una fase global. La condición C5 añade 12 compuertas
$\sqrt{X}$ (\texttt{SX}) y 12 compuertas $\sqrt{X}^\dagger$ (\texttt{SXdg}) respecto a C4,
con lo que la profundidad del circuito pasa de 65 a 67 capas, sin compuertas CNOT
adicionales. Al restringirse ambos mapas a interacciones de grado $|S| \le 2$, C5
no introduce interacciones de orden superior respecto a C4.
